\chapter{Finite propagation speed, \texorpdfstring{$L^2$}{L2} estimates, and bounded sectional curvature}
\source{cheeger-gromov-taylor:page-0005,cheeger-gromov-taylor:page-0006,cheeger-gromov-taylor:page-0007,cheeger-gromov-taylor:page-0008,cheeger-gromov-taylor:page-0009}

Let $M$ be complete. The estimates below use finite propagation for the wave
operator and, where stated, the curvature bound $\lvert K_M\rvert\leq K$.

Recall that $u=\cos s\sqrt{-\Delta}\,\delta_{x_2}$ satisfies the wave equation

\begin{equation}\tag{1.1}
\left(\frac{\partial^2}{\partial s^2}-\Delta\right)u=0,
\end{equation}

with initial data

\begin{equation}\tag{1.2}
u(0,x)=\delta_{x_2}, \qquad \frac{\partial}{\partial s}u(0,x)=0.
\end{equation}

Finite propagation for the wave equation gives

\begin{equation}\tag{1.3}
\supp \cos s\sqrt{-\Delta}\,\delta_{x_2}\subset \overline{B_{|s|}(x_2)},
\end{equation}

i.e., $\cos s\sqrt{-\Delta}$ has \emph{unit} propagation speed. This fact can be
brought to bear on the estimation of the kernels of other functions of
$\sqrt{-\Delta}$, by employing the formula, valid for an even function $f(\lambda)$:

\begin{equation}\tag{1.4}
f(\sqrt{-\Delta})=\frac{1}{2\pi}\int_{-\infty}^{\infty}\hat f(s)\cos s\sqrt{-\Delta}\,ds,
\end{equation}

where

\begin{equation}\tag{1.5}
\hat f(s)=\int_{-\infty}^{\infty}f(\lambda)\cos \lambda s\,d\lambda.
\end{equation}

Formula (1.4) is the even Fourier inversion formula for $\sqrt{-\Delta}$.

\begin{definition}[Finite-propagation class]
\label{def:finite-propagation-class}
\source{cheeger-gromov-taylor:page-0006}
\dcref{Definition 1.1}
\group{section-1}

We say
\begin{equation}\tag{1.6}
f \in \Tscr_1^k,
\end{equation}
provided that $f$ is even and, for each $I = (-\varepsilon, \varepsilon)$,
\begin{equation}\tag{1.7}
\left(\frac{d}{ds}\right)^l \hat{f}(s) \in L^1(\R \setminus I), \qquad 0 \leq l \leq k.
\end{equation}
\end{definition}

\begin{proposition}[Tail estimate for $f(\sqrt{-\Delta})$]
\label{prop:tail-estimate}
\source{cheeger-gromov-taylor:page-0006}
\dcref{Proposition 1.1}
\group{section-1}

Let $M^n$ be complete, $x_2\in M^n$, and let $B_R(x_2)$ be the radius-$R$ ball.
If $f\in\Tscr_1$, then
\begin{equation}\tag{1.8}
\left\| f(\sqrt{-\Delta})u \right\|_{M \setminus B_R(x_2)}
\leq \|u\| \frac{1}{\pi} \int_{R-r}^{\infty} |\hat{f}(s)|\, ds
\end{equation}
provided
\begin{equation}\tag{1.9}
\supp u \subset B_r(x_2), \qquad r < R.
\end{equation}
\end{proposition}

\begin{proof}
\source{cheeger-gromov-taylor:page-0006}
\sketch Using (1.4), finite propagation (1.3) and (1.9) give
\begin{equation}\tag{1.10}
f(\sqrt{-\Delta})u = \frac{1}{2\pi} \int_{-\infty}^{\infty} \hat{f}(s)\cos s\sqrt{-\Delta}\,u\,ds.
\end{equation}
\begin{equation}\tag{1.11}
\cos s\sqrt{-\Delta}\,u = 0 \quad \text{on } M \setminus B_{r+|s|}(x_2).
\end{equation}
Hence the left side of (1.8) is bounded by
\[
\left\| \frac{1}{\pi} \int_{R-r}^{\infty} \hat{f}(s)\cos s\sqrt{-\Delta}\,u\,ds \right\|
\leq \frac{1}{\pi} \int_{R-r}^{\infty} \|u\|\,|\hat{f}(s)|\,ds,
\]
since
\begin{equation}\tag{1.12}
\|\cos s\sqrt{-\Delta}\| \leq 1.
\end{equation}
\end{proof}

\begin{corollary}[Derivative tail estimate]
\label{cor:derivative-tail-estimate}
\source{cheeger-gromov-taylor:page-0006}
\dcref{Corollary 1.2}
\group{section-1}

Under the hypotheses of Proposition~\ref{prop:tail-estimate}, if $f\in\Tscr_1^{2k+2l}$, then
\begin{equation}\tag{1.13}
\left\| \Delta^k f(\sqrt{-\Delta}) \Delta^l u \right\|_{M \setminus B_R(x_2)}
\leq \|u\| \cdot \frac{1}{\pi} \int_{R-r}^{\infty} \left|\hat{f}^{(2k+2l)}(s)\right|\, ds.
\end{equation}
\end{corollary}

\begin{proof}
\source{cheeger-gromov-taylor:page-0006}
\sketch Differentiating (1.4) under the integral gives
\begin{equation}\tag{1.14}
\Delta^k f(\sqrt{-\Delta}) \Delta^l u
= \frac{1}{2\pi} \int_{-\infty}^{\infty} \hat{f}^{(2k+2l)}(s)\cos s\sqrt{-\Delta}\,u\,ds,
\end{equation}
and the proof of Proposition~\ref{prop:tail-estimate} applies.
\end{proof}

Assume $\lvert K_M\rvert\leq K$. Choose
\begin{equation}\tag{1.15}
 r_1\le \min\bigl(K^{-1/2},\ii(x_1)\bigr),\qquad
 r_2\le \min\bigl(K^{-1/2},\ii(x_2)\bigr).
\end{equation}
For the local estimate, first assume
\begin{equation}\tag{1.16}
 r_j=1,\qquad \lvert K_M\rvert \leq 1,
\end{equation}
and that (1.15) holds. Recall that if $g(r)$ is any function of $r$ alone in a
normal ball $B_{r_0}(x_0)$, we have
\begin{equation}\tag{1.17}
 \Delta g = g'' + \frac{A'}{A}g'.
\end{equation}
Here $A(r,\theta)$ is the polar area density. Comparison gives, if
$\lvert K_M\rvert\leq1$, on $B_1(x_0)$
\begin{equation}\tag{1.18}
 0<\frac{c_1(n)}{r}<\frac{A'}{A}<\frac{c_2(n)}{r}.
\end{equation}
Let $\phi$ be a smooth function supported on $[0,1/2]$ with
$\phi\,|_{[0,1/4]}\equiv 1$, $|\phi'|<5$, and set $\phi_{\eps}=\phi(r/\eps)$.
Let
\[
 P=\phi_{\eps}(r)\,\frac{r^{2-n}}{\alpha(n-1)},\qquad
\left(P-\phi_{\eps}\frac{\log r}{2\pi},\ n=2\right).
\]
Finally let $\Delta_x(P(x_0,x))=Q$ denote $\Delta_x$ applied pointwise to $P$.
Then on $B_1(x_0)$,
\begin{equation}\tag{1.19}
 |Q(r)|\le C_{\eps}r^{2-n}.
\end{equation}
For smooth $g$ on $B_1(x_0)$,
\begin{equation}\tag{1.20}
 \int_{B_1(x_0)} P(x_0,x)\Delta g(x)\,dx = g(x_0)-\int_{B_1(x_0)} Q(x)g(x)\,dx.
\end{equation}
By the above mentioned comparison, the metric on $B_1(x_0)$ is bounded
uniformly above and below by a constant times the Euclidean metric in normal
coordinates. Fix $\eps$ sufficiently small so that the $N^2$ compositions below
are defined and apply the standard iteration argument (see e.g. \cite{cgt-de-rham}) to (1.20).
Thus if we write (1.20) as
\begin{equation}\tag{1.21}
 P\Delta = I - Q,
\end{equation}
and set
\begin{equation}\tag{1.22}
 \Pscr = (I + \Qscr + \cdots + \Qscr^{N-1})P,\qquad
 N=\left[\frac{n}{4}\right]+1,
\end{equation}
\[
\Qscr = Q^N,
\]
then it follows easily that
\begin{equation}\tag{1.23}
 \Pscr^N(\Delta^N g) + \Pscr^{N-1}\Qscr(\Delta^{N-1}g) + \cdots + \Pscr\Qscr(\Delta g) + \Qscr g = g.
\end{equation}
The kernels on the left hand are continuous and their sup norms are bounded,
since the metric is boundedly related to the Euclidean metric in normal
coordinates, and (1.19) holds. Thus we get
\begin{equation}\tag{1.24}
|g(x_0)| \leq c(n)\bigl(\|g\|_{B_{i(x_0)}} + \|\Delta g\|_{B_{i(x_0)}} + \cdots + \|\Delta^N g\|_{B_{i(x_0)}}\bigr),
\qquad
N=\left[\frac{n}{4}\right]+1.
\end{equation}

Scaling the metric gives

\begin{proposition}[Local elliptic estimate]
\label{prop:local-elliptic-estimate}
\source{cheeger-gromov-taylor:page-0008}
\dcref{Proposition 1.3}
\group{section-1}

Let $\lvert K_M\rvert\leq K$.

\begin{equation}\tag{1.25}
r \leq \min\bigl(|K|^{-1/2}, \ii(x_0)\bigr),\qquad
N=\left[\frac{n}{4}\right]+1,
\end{equation}

Then

\begin{equation}\tag{1.26}
|g(x_0)| \leq c(n)\bigl(r^{-n/2}\|g\|_{B_r(x_0)} + \cdots + r^{2N-n/2}\|\Delta^N g\|_{B_r(x_0)}\bigr).
\end{equation}
\end{proposition}

In Proposition~\ref{prop:refined-injectivity-angle-lower-estimates} we will give a lower estimate for $\ii(x)$ which will be
incorporated into Theorem~\ref{thm:kernel-estimates-bounded-curvature} below. First we introduce some notation.
Let $\Vr(p)$ denote the volume of $B_r(p)$ in $M^n$, and $\VHr$ the volume of
$B_r(q)\subset M_H^n$, where $M_H^n$ is the simply connected $n$-dimensional
space of curvature $H$. Thus

\begin{equation}\tag{1.27}
\VHr = \alpha(n-1)\int_0^r \QH(s)\,ds,
\end{equation}

where

\begin{equation}\tag{1.28}
\QH(s)=
\begin{cases}
\left(\dfrac{\sin \sqrt{H}\,s}{\sqrt{H}}\right)^{n-1}, & H>0,\\[1.2ex]
s^{n-1}, & H=0,\\[1.2ex]
\left(\dfrac{\sinh \sqrt{|H|}\,s}{\sqrt{|H|}}\right)^{n-1}, & H<0,
\end{cases}
\end{equation}

and $\alpha(n-1)=V(S^{n-1}_1)$.

\begin{theorem}[Kernel estimates for bounded sectional curvature]
\label{thm:kernel-estimates-bounded-curvature}
\source{cheeger-gromov-taylor:page-0009}
\dcref{Theorem 1.4}
\group{section-1}

Let $M^n$ be complete with $H \leq \KM \leq K$. Fix $p$, $x_1$, $x_2 \in M^n$,
$r_0 > 0$. Set $\rho_j = \rho(p,x_j)$, $j=1,2$, and $d = \rho(x_1,x_2)$.
Let $4N > n$ and assume $\hat f^{(k)} \in L^1(\R)$, $0 \le k \le 4N$.
Fix $r$, $r_0$, $s$ with $r_0 + 2s \leq \pi/\sqrt{K}$ and $r_0 \leq \pi/4\sqrt{K}$.

\textup{(i)} The kernel $k_f$ of $f(\sqrt{-\Delta})$ is continuous on
$\{(x_1,x_2) \in M \times M:\ x_1 \ne x_2\}$ and satisfies

\begin{equation}\tag{1.29}
|k_f(x_1,x_2)| \leq \frac{1}{\pi} c^2(n) \sum_{0 \le i,j \le N}
\bar r_1^{2i-n/2}\bar r_2^{2j-n/2} \int_0^\infty \bigl|\hat f^{(2i+2j)}(s)\bigr|\,ds,
\end{equation}

where $c(n)$ is the constant in (1.26), and

\begin{equation}\tag{1.30}
\bar r_j = \min\!\left(|H|^{-1/2}, \frac{r_0}{2}\,\frac{1}{1 + \left(V^{H}_{r_0+s}/V_r(p)\right)\left(V^{H}_{\rho_j+r}/V^{H}_s\right)}\right).
\end{equation}

\textup{(ii)} If only $f \in \mathcal{F}^{4N}_1$, then for $r_j < \bar r_j$, $j=1,2$,
and $r_1 + r_2 < d$,

\begin{equation}\tag{1.31}
|k_f(x_1,x_2)| \leq \frac{2}{\pi} c^2(n) \sum_{0 \le i,j \le N}
r_1^{2i-n/2}r_2^{2j-n/2} \int_{d-r_1-r_2}^\infty \bigl|\hat f^{(2i+2j)}(s)\bigr|\,ds.
\end{equation}
\end{theorem}

\begin{proof}
\source{cheeger-gromov-taylor:page-0009}
\uses{prop:local-elliptic-estimate,cor:derivative-tail-estimate}
\sketch The arguments for (1.29) and (1.31) are the same; consider (1.31).
Let $u \in L^2$ be supported on
$B_{r_2}(x_2)$. We apply (1.26) of Proposition~\ref{prop:local-elliptic-estimate} with
$g = \Delta_2^j k_f(x_1,x_2)(u(x_2))$, $x_0 = x_1$; by (1.30) and the estimate
for $i(x_1)$ of Theorem~\ref{thm:injectivity-radius-relative-volume}, the hypothesis (1.25) is satisfied. If we combine
the resulting estimate with (1.13) of Corollary~\ref{cor:derivative-tail-estimate}, we get

\begin{equation}\tag{1.32}
\bigl|\Delta_2^j k_f(x_1,x_2)(u)\bigr| \leq \frac{c(n)}{\pi}\sum_{i=0}^N r_1^{2i-n/2}
\bigl\|\Delta_1^i \Delta_2^j k_f(u)\bigr\|_{B_{r_1}(x_1)}
\end{equation}
\[
\leq c(n)\|u\|_{B_{r_2}(x_2)} \sum_{i=0}^N r_1^{2i-n/2}
\int_{d-r_1-r_2}^\infty \bigl|\hat f^{(2i+2j)}(s)\bigr|\,ds.
\]

Since $u$ is arbitrary, for all $x_1$ we have

\begin{equation}\tag{1.33}
\bigl\|\Delta_2^j k_f(x_1,x_2)\bigr\|_{B_{r_2}(x_2)}
\leq c(n)\sum_{i=0}^N r_1^{2i-n/2} \int_{d-r_1-r_2}^\infty \bigl|\hat f^{(2i+2j)}(s)\bigr|\,ds.
\end{equation}

Applying (1.26) on $B_{r_2}(x_2)$ proves (1.31).
\end{proof}

\begin{remark}
\label{rem:sectional-curvature-kernel-estimate}
\source{cheeger-gromov-taylor:page-0009}
\dcref{Remark after Theorem 1.4}
\group{section-1}

The estimate for $i(x_j)$ in Theorem~\ref{thm:kernel-estimates-bounded-curvature} can be sharpened when $\rho_j$ is
sufficiently large, as in Theorem~\ref{thm:injectivity-radius-relative-volume}(ii). For compact $M^n$, known injectivity
radius bounds may be substituted.
\end{remark}

Theorem~\ref{thm:kernel-estimates-bounded-curvature} extends to $\partial M^n\ne\emptyset$ using a suitable
generalization of Proposition~\ref{prop:refined-injectivity-angle-lower-estimates}.
